\documentclass[12pt,letterpaper]{article}

% --- Packages -----------------------------------------------------------------
\usepackage[utf8]{inputenc}
\usepackage[T1]{fontenc}
\usepackage{amsmath,amssymb,amsthm}
\usepackage{mathptmx}
\usepackage[margin=1in]{geometry}
\usepackage{setspace}
\usepackage{graphicx}
\usepackage{booktabs}
\usepackage{natbib}
\usepackage{hyperref}
\usepackage{xcolor}
\usepackage{enumitem}

\newtheorem{assumption}{Assumption}
\newtheorem{proposition}{Proposition}
\newtheorem{definition}{Definition}
\newtheorem{lemma}{Lemma}

\hypersetup{colorlinks=true,linkcolor=blue,citecolor=blue,urlcolor=blue}

\title{A Structural Econometric Model of AI Demand, Generator Investment,
and the Ancillary Services Market}
\author{Aruna Balasubramanian, Niranjan Balasubramanian, and Dana Golden}
\date{May 2026}


\begin{document}

\maketitle

\begin{abstract}
\noindent We develop a structural econometric model of the interaction between
hyperscaler data center demand and the U.S.\ electricity market, integrating
energy, capacity, and ancillary services. The model has three time scales: a
short run in which a network energy and ancillary services market clears, a
medium run in which generators enter, exit, contract via PPAs, and bid into
capacity markets, and a long run in which hyperscalers choose data center
capacity and procurement portfolios across an exogenously fixed but
counterfactually adjustable set of locations. Hyperscaler payoffs are driven
by an investor-capital contest with decreasing marginal returns to compute,
following a now-standard formulation in the AI-economics literature.
Energy and ancillary-services cost curves are imported exogenously from
ISO-published data and futures prices, while the impact of AI workloads on
ancillary-services demand is identified from a stacked difference-in-differences
design using model release dates as the treatment timing.
The model is connected to a coupled microgrid--micro-data-center digital twin
that produces experimental measurements of how data center design parameters
(quantization, batch scheduling, storage configuration, behind-the-meter
generation, load-shaping) translate into voltage, harmonic, ramp-rate, and
inertia signatures at the point of grid connection. These mappings provide
the link from data-center-level technical choices to grid-level ancillary
services demand. We use the estimated model to evaluate counterfactuals
including geographic redistribution of training versus inference workloads,
technical changes at the model level, bring-your-own-clean-generation mandates,
ancillary services market design, and 24/7 carbon-free matching.
\end{abstract}

\newpage

\section{Introduction}
\label{sec:intro}

The growth of artificial intelligence workloads is the most consequential
shift in U.S.\ electricity demand of the past decade. Forecasts for
data-center load growth are sharply at variance with the historical experience
of essentially flat aggregate U.S.\ electricity demand, and the heterogeneity
of impact is large: load growth is geographically concentrated, exhibits
load profiles unlike traditional industrial or residential customers, and is
operated by a small number of firms (``hyperscalers'') whose internal design
choices have macroscopic grid consequences \citep{shehabi2024,epri2024,
gridstrategies2025}.

Two facts make the policy question difficult. First, the technology of the
load is unusual. AI training is large, sustained, and synchronized across
GPU clusters; AI inference is bursty, latency-sensitive, and increasingly
geographically dispersed. Both interact with the grid through power electronics
whose ramp-rate, harmonic, and reactive-power characteristics differ from
conventional industrial customers in ways that the existing power-quality
literature has only begun to document \citep{LaughnerEtAl_2024_THDonUSGrid,
WhiskerLabs_THD_2024}. Second, the policy levers are heterogeneous. Hyperscalers
do not simply purchase electricity. They choose locations, contract structures
(merchant, REC, PPA, colocation), backup configurations (diesel, battery,
hydrogen), and---increasingly---procurement of behind-the-meter generation.
Each lever has different implications for emissions, reliability, ancillary
services, and ratepayer cost.

Existing reduced-form work, including our own prior analysis using stacked
difference-in-differences around model release dates, has established that
AI activity raises wholesale prices, increases fossil generation, and
degrades power quality in zones with elevated data-center exposure. The
reduced-form approach has two limitations that motivate the present paper.
First, it cannot speak to mechanism: the same reduced-form effect could be
generated by load growth, by procurement choices that displace clean
generation, or by power-electronics signatures that increase ancillary
services costs. Second, it cannot speak to counterfactuals that change the
choice environment: a 24/7 matching mandate, a colocation requirement, or
a change in PJM capacity market design will reallocate equilibrium
quantities in ways that a reduced-form estimate cannot evaluate.

This paper develops a structural model designed to address both problems. The
model has three central features. \emph{First}, the energy market and most
of the ancillary services market are solved exogenously, importing
ISO-published cost curves and futures prices as primitives. This is a
deliberate tractability choice: it sacrifices endogenous wholesale-price
formation in order to make the demand-side hyperscaler problem and the
generator entry margin computationally feasible. \emph{Second}, the impact of
AI workloads on ancillary services demand is identified empirically rather
than imposed by assumption, using a stacked difference-in-differences design
in the spirit of our prior reduced-form work. \emph{Third}, the mapping from
data-center-level design parameters to grid-level signatures is calibrated
using a coupled microgrid--micro-data-center digital twin under
development at Stony Brook University's CEWIT/AERTC. The digital twin
provides the experimental variation in model design (quantization,
distillation, batch scheduling), interconnection design (storage,
behind-the-meter generation), and load shape that is impossible to recover
from observational data alone.

The model is structural in the sense that it recovers preference and cost
primitives---hyperscaler reputational sensitivity, financing-risk cost-of-capital,
contest concentration, generator entry costs by fuel type---from observed
choices using a moment-inequalities estimator on the long-run problem and
two-stage methods on the short-run and medium-run problems. With those
primitives in hand we solve counterfactual Markov-perfect equilibria for
each of the policy scenarios listed in the design.

\paragraph{Contribution.} The paper makes three contributions. (i) It
integrates a hyperscaler combinatorial discrete-choice problem with an
operations-research-style ancillary services market into a single estimable
framework. To our knowledge no prior work has done this. (ii) It provides
the first empirical estimates of AI workload effects on the ancillary
services market specifically (as distinct from energy or capacity markets),
which is the segment of the wholesale market in which the most distinctive
features of AI demand (sub-second ramps, harmonic content, reactive power)
matter. (iii) It develops the experimental-econometric integration: digital
twin measurements identify the scale-free \emph{derivative} of grid signatures
with respect to technical design parameters---variation no observational
dataset contains---while observational data supply the \emph{level}, and the
two are reconciled at the zone level as an overidentifying test rather than a
calibration.

\paragraph{Road map.} Section~\ref{sec:lit} situates the paper in five
literatures (structural models of electricity markets, dynamic oligopoly with
lumpy capacity, combinatorial discrete choice and spatial sorting, the
empirical economics of AI and the grid, and modern difference-in-differences).
Section~\ref{sec:model} develops the network and the long-run hyperscaler and
generator problems; Section~\ref{sec:shortrun} develops the short-run
energy, REC, capacity, and ancillary-services markets and the equilibrium
concept. Section~\ref{sec:data} describes the data. Section~\ref{sec:estimation}
develops the three-stage estimation strategy. Section~\ref{sec:experiments}
describes the physical experiments and their integration with the
econometric model. Sections~\ref{sec:results}--\ref{sec:cf} present results,
the decomposition of AI's system impacts, and counterfactuals.

\section{Literature Review}
\label{sec:lit}

The paper sits at the intersection of five literatures, and its contribution
is best understood as the seam it stitches between them.

\paragraph{Structural models of wholesale electricity markets.} A long
tradition recovers cost and conduct primitives from bid and dispatch data.
\citet{wolfram1999measuring} measures market power in the British pool from the
gap between bids and marginal cost; \citet{cicala2022imperfect} contrasts the
inefficiency of imperfect competition with that of imperfect regulation across
U.S.\ generation; and \citet{mercadal2016electricity} studies how financial
participants reshape arbitrage and price formation in organized markets. These
papers treat the supply stack and the price-formation mechanism as the
objects of interest. We depart from them on exactly this margin: rather than
endogenize wholesale price formation, we import the energy and capacity price
processes as data (\S\ref{sec:shortrun}) and reinvest the freed degrees of
freedom in a richer demand side. The cost is muted within-period price
feedback; the benefit is that the hyperscaler's combinatorial procurement
problem becomes estimable. The one wholesale segment we do treat as a
behavioral object---ancillary services---is precisely the one the prior
structural literature has left aside.

\paragraph{Dynamic oligopoly with lumpy capacity.} The machinery for the
long-run game comes from the literature on dynamic investment under sunk,
discrete capacity \citep{gowrisankaran2012capacity}, and from the
revealed-preference, moment-inequality approach to estimating such games
without solving for a unique equilibrium \citep{holmes2011walmart}. We borrow
the existence argument for dynamic discrete games with capacity as a state and
the inferential posture---partial identification rather than point
identification---that follows when uniqueness is not guaranteed
(\S\ref{sec:est_long}). The novelty relative to this literature is the
two-sided structure: hyperscaler entry and generator entry are coupled through
the contract market, so a hyperscaler's procurement choice is also a
supply-side instrument.

\paragraph{Combinatorial discrete choice and spatial sorting.} Hyperscalers
choose a \emph{set} of locations and procurement modes, not a single one, and
the value of a location depends on the others chosen through the firm's
aggregate compute. This is a combinatorial discrete-choice problem of the kind
formalized by \citet{arkolakis2023combinatorial}, whose monotonicity-based
location-pruning algorithm we adapt in \S\ref{sec:hyper_lr}. The adaptation is
not mechanical: the contest payoff couples locations through a scalar
aggregate, and the resource constraints couple firms, so the conditions under
which pruning is valid must be stated with care (\S\ref{sec:est_long},
Appendix~B).

\paragraph{The empirical economics of AI and the grid.} A fast-growing applied
literature documents the scale of data-center load growth
\citep{shehabi2024, epri2024, gridstrategies2025} and its energy and emissions
footprint \citep{deVries2023, Thompson2024_AI_Datacenters_Energy_Dynamics}, and
an emerging strand connects data-center proximity to measurable power-quality
degradation \citep{LaughnerEtAl_2024_THDonUSGrid, WhiskerLabs_THD_2024}. This
work, including our own prior reduced-form analysis, establishes that AI
activity moves wholesale outcomes; it does not isolate the mechanism or
evaluate counterfactual choice environments. The present paper takes the
reduced-form findings as a target the structural model must reproduce
(\S\ref{sec:est_short}) and then asks the mechanism and policy questions the
reduced-form approach cannot.

\paragraph{Modern difference-in-differences and partial-identification
inference.} The first-stage identification of the ancillary-services impact
uses the recent event-study toolkit for staggered treatment
\citep{callaway2021difference, borusyak2024revisiting}, and the long-run
estimator uses moment-inequality inference \citep{andrews2010inference}. These
are methods we apply rather than extend; we flag them here because the
credibility of the structural estimates rests on the credibility of the
reduced-form moment they are asked to match.

\paragraph{Renewable matching and the timing wedge.} Finally, the emissions
accounting and the 24/7 carbon-free-electricity counterfactual
(\S\ref{sec:cf}) draw on the literature on the wedge between annual and hourly
renewable matching. We position the granularity standard $h$ as a continuous
policy lever spanning that literature's two poles.\footnote{The 24/7 CFE and
hourly-matching references are to be added from the companion bibliography.}


\section{Model}
\label{sec:model}

The model is developed across this section and the next. The present section
fixes the network (\S\ref{sec:network}) and the long-run problems of the two
endogenous agent classes, hyperscalers and generators (\S\ref{sec:longrun}).
Section~\ref{sec:shortrun} then develops the short-run energy, capacity, REC,
and ancillary-services markets---most solved using exogenously imported cost
curves---and closes with the timing of moves and the equilibrium concept
(\S\ref{sec:timing}). The ancillary-services market (\S\ref{sec:aux_market}) is
the one short-run market treated as a behavioral object rather than imported,
and is the locus of the paper's identification.

\paragraph{Overview of endogeneity and exogeneity.} A central modeling decision
is which markets are solved endogenously and which are imported as data; this
paragraph states the split and its cost once, and later sections refer back to
it rather than re-arguing the point.
Endogenous: hyperscaler capacity, procurement-mode, and storage choices;
generator capacity choices by fuel type and location; the PPA contracting
margin; the entry margin for new generators; the impact of AI workloads on
ancillary services demand. Exogenous: short-run energy dispatch prices (taken
from ISO LMP data and futures prices); short-run ancillary services prices
(taken from ISO ancillary services market data); capacity market clearing
prices (taken from PJM Base Residual Auction results and analogous data from
other ISOs). This is a deliberate tractability decision: it isolates the
mechanisms that the design proposes to study (hyperscaler demand-side levers
and the supply-side investment response) without requiring solution of a
full unit-commitment optimal-power-flow problem at each iteration of the
estimation routine. The cost of this simplification is that price feedback
from the long-run capacity decisions back to short-run market clearing is
muted; we address this by re-running the short-run market with a perturbed
supply stack in counterfactuals (\S\ref{sec:cf}) so that capacity changes do
shift prices, even though the short-run prices entering the estimation are
treated as given.
generator capacity choices by fuel type and location; the PPA contracting
margin; the entry margin for new generators; the impact of AI workloads on
ancillary services demand. Exogenous: short-run energy dispatch prices (taken
from ISO LMP data and futures prices); short-run ancillary services prices
(taken from ISO ancillary services market data); capacity market clearing
prices (taken from PJM Base Residual Auction results and analogous data from
other ISOs). This is a deliberate tractability decision: it isolates the
mechanisms that the design proposes to study (hyperscaler demand-side levers
and the supply-side investment response) without requiring solution of a
full unit-commitment optimal-power-flow problem at each iteration of the
estimation routine. The cost of this simplification is that price feedback
from the long-run capacity decisions back to short-run market clearing is
muted; we address this by re-running the short-run market with a perturbed
supply stack in counterfactuals (\S\ref{sec:cf}) so that capacity changes do
shift prices, even though the short-run prices entering the estimation are
treated as given.

\subsection{Network}
\label{sec:network}

\subsubsection{Network Primitives}

\paragraph{Echelons.} Let $\mathcal{E} = \{1, 2, \dots, E\}$ index echelons,
ordered so that bulk power transit proceeds from lower to higher indices.
We use $E = 3$ throughout as the canonical case (upstream / transit /
downstream), but the formulation extends to arbitrary $E$.

\paragraph{Zones.} Within echelon $e$, let $\mathcal{Z}_e = \{1, \dots, Z_e\}$
be the set of zones, where each zone is a pricing area (e.g., an ISO
sub-region or a balancing authority subarea). The full set of zones is
$\mathcal{Z} = \bigcup_{e \in \mathcal{E}} \mathcal{Z}_e$, and we write
$z \in \mathcal{Z}_e$ or equivalently $e(z) = e$ for the echelon of zone $z$.

\paragraph{Directed inter-echelon arcs.} The transmission network is a
directed graph $G = (\mathcal{Z}, \mathcal{A})$ where every arc
$a = (i, j) \in \mathcal{A}$ satisfies $e(j) = e(i) + 1$. That is, arcs
connect only adjacent echelons in the forward direction. This is the key
structural departure from a bipartite generator--load graph and from an
unrestricted nodal model: arcs are constrained to respect the echelon
ordering, ruling out backflow from downstream to upstream zones.\footnote{The
restriction is not a physical claim about power flow---real AC systems are
bidirectional and governed by Kirchhoff's laws---but a modeling choice that
captures the \emph{net} commercial direction of energy in markets where
generation and load are geographically segregated. For markets where flows
reverse seasonally or diurnally, echelon ordering can be relaxed to a
partial order, or two parallel arc sets $(\mathcal{A}^+, \mathcal{A}^-)$
can be introduced.}

\paragraph{Intra-echelon arcs (optional).} For richer applications, we admit
intra-echelon arcs $\mathcal{A}^{\parallel}_e \subseteq \mathcal{Z}_e \times
\mathcal{Z}_e$ representing transfers between zones at the same level
(e.g., interties between neighboring balancing areas). When
$\mathcal{A}^{\parallel}_e = \emptyset$ for all $e$, the network is purely
cascading.

\paragraph{Transfer capacities.} Each arc $a \in \mathcal{A} \cup
\bigcup_e \mathcal{A}^{\parallel}_e$ carries a transfer limit $\bar T_{a,t}$
in period $t$. We treat $\bar T_{a,t}$ as exogenous in the baseline (taken
from FERC Form~715 and ISO topology disclosures); a counterfactual lever
allows transmission expansion at locations of policy interest.

\paragraph{Generators and loads at the zone.} Each zone $z$ hosts a set of
generators $\mathcal{G}_z$ indexed by fuel type $l \in \mathcal{L}^{\text{fuel}}
= \{\text{coal, CCGT, CT, nuclear, solar, wind, hydro, BESS, other}\}$ and
a set of demand sources including a hyperscaler-attributable component
$D^M_{z,t}$ and a residual non-hyperscaler component $D^O_{z,t}$. The latter
is exogenous in the model (taken from EIA Form~930 hourly load data minus
estimated data-center load); the former is the central object of the
hyperscaler problem developed below.

\paragraph{Why the echelon structure matters here.} The echelon network is
appropriate for the AI-grid setting because hyperscaler data centers cluster
in a small number of downstream zones (Northern Virginia, central Ohio,
central Texas, Phoenix, Quincy/Boardman) while incremental generation
investments---particularly behind-the-meter and colocated generation---can
in principle be sited at either downstream or upstream zones depending on
the PPA structure. The echelon ordering captures the fact that PPAs sourcing
upstream-sited renewables for downstream data center load require transmission
allocation across the directed arc set, and that the timing wedge between
renewable generation and AI load is in part a function of transmission
distance.

\subsection{Long-run}
\label{sec:longrun}

The long-run is the time scale at which capital is committed: hyperscalers
build data centers, generators build or retire capacity, and PPAs and
colocation contracts are signed. We model the long-run as a dynamic game
with annual periods indexed by $t = 0, 1, 2, \dots$ and discount factor
$\beta \in (0,1)$. The state at the start of $t$ is

\begin{equation}
\label{eq:lr_state}
\mathbf{s}_t \;=\; \big(\,
\underbrace{\{K_{f,z,t}^\sigma\}_{f,z,\sigma}}_{\text{hyperscaler stocks}},\,
\underbrace{\{O^{\max}_{l,k,z,t}\}_{l,k,z}}_{\text{generator capacities}},\,
\underbrace{\{C^{\text{contract}}_{f,k,t}\}}_{\text{contract book}},\,
\underbrace{\{H_{z,t}, Z_{z,t}\}_{z}}_{\text{resource states}}\,
\big),
\end{equation}

where $\sigma \in \mathcal{I} = \{\text{merch, REC, PPA, Colo}\}$ is the
procurement mode, $K_{f,z,t}^\sigma$ is firm $f$'s installed data center
capacity (MW of IT load) at zone $z$ under mode $\sigma$, $O^{\max}_{l,k,z,t}$
is generator firm $k$'s capacity of fuel type $l$ at zone $z$,
$C^{\text{contract}}_{f,k,t}$ is the outstanding contract book between
hyperscaler $f$ and generator $k$, $H_{z,t}$ is local interconnection
headroom, and $Z_{z,t}$ is the local clean-MW pool. Locations are
\emph{exogenously given} in the baseline but can be perturbed in
counterfactuals: the policy lever ``relocate inference to upstream zones''
is implemented by altering the set of admissible $z$ for each $f$.

\subsubsection{Hyperscaler Problem}
\label{sec:hyper_lr}

A finite set $\mathcal{F} = \{1,\dots,F\}$ of hyperscalers makes capacity,
procurement, and storage choices at each $t$. Each $f$ is characterized by
two structural parameters: a productivity shifter $\theta_f$ mapping compute
into model quality, and a reputational sensitivity $\gamma_f$ that determines
the cost-of-capital penalty for emissions. Both are private types drawn from
distributions $F_\theta, F_\gamma$ to be recovered in estimation.

\paragraph{Action space.} In each period $t$, firm $f$ at each admissible
zone $z$ chooses
\begin{equation}
\label{eq:hyper_action}
a_{f,z,t} \;=\; \big(\,
\Delta K_{f,z,t},\; \sigma_{f,z,t},\; \Delta S_{f,z,t},\; B_{f,z,t}\,
\big),
\end{equation}
where $\Delta K_{f,z,t} \in \{0\} \cup \mathcal{K}$ is the increment to
data center capacity at $z$ (drawn from a finite grid $\mathcal{K} =
\{k_1, \dots, k_M\}$ of standardized building-block sizes),
$\sigma_{f,z,t} \in \mathcal{I}$ is the procurement mode chosen for the
increment, $\Delta S_{f,z,t} \in \{0\} \cup \mathcal{S}$ is the increment
to behind-the-meter storage capacity (MWh), and $B_{f,z,t} \in \mathcal{B}^{\mathrm{bk}} =
\{\text{none, diesel, BESS, hybrid}\}$ is the backup configuration. (We reserve
$\mathcal{B}$ for the set of short-run load blocks introduced in
\S\ref{sec:shortrun}; the backup-configuration set carries the superscript
$\mathrm{bk}$ throughout to avoid collision with the block index $b$.) The
discretization of $(\Delta K, \Delta S)$ is necessary for the moment-inequalities
estimator and reflects engineering modularity (data centers are built in
standardized blocks).

\paragraph{State transition.} Capacity accumulates with depreciation rate
$\delta_K$:
\begin{equation}
\label{eq:state_transition}
K_{f,z,t+1}^\sigma = (1 - \delta_K) K_{f,z,t}^\sigma + \Delta K_{f,z,t}
\cdot \mathbf{1}\{\sigma_{f,z,t} = \sigma\}, \qquad
S_{f,z,t+1} = (1 - \delta_S) S_{f,z,t} + \Delta S_{f,z,t}.
\end{equation}
Capacity is sunk: $\Delta K_{f,z,t} \geq 0$.

\paragraph{Compute and quality.} Total compute available to $f$ at $t$ is
\begin{equation}
\label{eq:compute}
C_{f,t} \;=\; \sum_{z \in \mathcal{Z}_f}\sum_{\sigma \in \mathcal{I}}
\kappa_z \cdot K_{f,z,t}^\sigma,
\end{equation}
where $\mathcal{Z}_f$ is firm $f$'s admissible location set and $\kappa_z$
is a zone-specific efficiency factor capturing PUE, climate-driven cooling
overhead, and chip generation availability. Model quality is a stochastic
increasing function of compute and the firm's productivity type:
\begin{equation}
\label{eq:quality}
q_{f,t} \;=\; \theta_f \cdot C_{f,t}^{\alpha_f} \cdot \varepsilon_{f,t},
\qquad \alpha_f \in (0,1],\; \varepsilon_{f,t} \sim F_\varepsilon \text{ i.i.d.}
\end{equation}
The exponent $\alpha_f < 1$ captures firm-specific decreasing marginal
returns to compute: firms whose models are bottlenecked by data quality or
research personnel have $\alpha_f$ well below one, generating compounding
disadvantage from over-investment in compute alone.

\paragraph{Investor capital contest.} Investors allocate a fixed pool of
capital $V_t$ across hyperscalers according to a Tullock contest over
expected quality:
\begin{equation}
\label{eq:contest}
v_{f,t} \;=\; V_t \cdot
\frac{\mathbb{E}_{t-1}[q_{f,t}^\eta]}{\sum_{f' \in \mathcal{F}} \mathbb{E}_{t-1}[q_{f',t}^\eta]},
\qquad \eta > 0.
\end{equation}
Investors' expectations are formed under their information set at $t-1$,
which we operationalize empirically using observed compute build-out as a
public signal. The exponent $\eta$ controls how sharply investment
concentrates on the perceived leader: $\eta \to \infty$ is winner-take-all,
$\eta \to 0$ is an equal split.

\paragraph{Hyperscaler per-period payoff.} Firm $f$'s flow payoff is
investor capital minus total cost:
\begin{equation}
\label{eq:hyper_flow}
\pi_{f,t} \;=\; v_{f,t} \;-\; C^{\mathrm{tot}}_{f,t} \;-\; C^{\mathrm{out}}_{f,t}
\;-\; \Gamma(\gamma_f, E_{f,t}) \cdot V_t,
\end{equation}
where $C^{\mathrm{tot}}_{f,t}$ is per-period operating and capital cost
detailed below, $C^{\mathrm{out}}_{f,t}$ is expected outage cost, and
$\Gamma(\gamma_f, E_{f,t})$ is a reputational damage function from emissions
$E_{f,t}$, scaled by the contestable investment pool $V_t$. The total cost is
\begin{equation}
\label{eq:hyper_cost}
\begin{aligned}
C^{\mathrm{tot}}_{f,t} \;=\; \sum_{z \in \mathcal{Z}_f}\Big[
& \underbrace{c^{\mathrm{land}}_z K_{f,z,t} + c^{\mathrm{chip}}_t K_{f,z,t}}_{\text{capital, annualized}}
+ \underbrace{c^{\mathrm{ene}}_{z,\sigma,t} \cdot u_{f,z,t}}_{\text{energy}}
+ \underbrace{c^{\mathrm{aux}}_{z,t} \cdot u^{\mathrm{aux}}_{f,z,t}}_{\text{ancillary services}} \\
& + \underbrace{c^{\mathrm{stor}} S_{f,z,t}}_{\text{storage capital}}
+ \underbrace{\Phi_\sigma(\Delta K_{f,z,t}) + \Phi_S(\Delta S_{f,z,t})}_{\text{procurement/storage setup}}
\Big].
\end{aligned}
\end{equation}
The four cost components that warrant comment are (i) the energy cost
$c^{\mathrm{ene}}_{z,\sigma,t}$, which is procurement-mode-specific and is
derived below in \S\ref{sec:proc_costs}; (ii) the ancillary-services cost
$c^{\mathrm{aux}}_{z,t}$, which is allocated to the hyperscaler in
proportion to its contribution to system ancillary-services demand (derived
in \S\ref{sec:aux_market}); (iii) the storage capital cost $c^{\mathrm{stor}}$,
which is paid per MWh of installed behind-the-meter storage capacity; and
(iv) the setup cost $\Phi_\sigma$, which captures one-time procurement
costs (PPA legal origination, colocation construction, interconnection
studies).

\paragraph{Outage cost.} Outage cost depends on the joint failure
probability of contracted and backup supply:
\begin{equation}
\label{eq:outage_cost}
C^{\mathrm{out}}_{f,t} \;=\; \sum_{z \in \mathcal{Z}_f}
P^{\mathrm{out}}(\sigma_{f,z,t}, B_{f,z,t}, S_{f,z,t}) \cdot
\bar L^{\mathrm{out}}_{f,z} \cdot K_{f,z,t},
\end{equation}
where $\bar L^{\mathrm{out}}_{f,z}$ is the value of lost compute per MW per
unit time at zone $z$ for firm $f$, and $P^{\mathrm{out}}(\cdot)$ is the
probability of unmet load. Storage and backup configurations reduce
$P^{\mathrm{out}}$; their values to the hyperscaler operate through this
term. Critically, storage also has grid-side effects on power quality and
ancillary services demand, which appear in the short-run market clearing
in \S\ref{sec:shortrun}; this dual role is the central reason storage is a
distinct choice variable rather than a sub-component of backup.

\paragraph{Emissions accounting.} Emissions attributable to $f$ at $t$ are
the sum across zones of grid-drawn energy times the marginal carbon
intensity (MCI) of the zone, net of REC retirements weighted by the
granularity standard. The granularity standard $h \in [0,1]$ is a
counterfactual lever (with $h=0$ annual matching, $h=1$ hourly matching).
Specifically, with $G^M_{f,z,b,t}$ denoting grid draw in load block $b$
(see \S\ref{sec:proc_costs}) and $R^{M,G}_{f,z,b,t}$ granular REC retirements,
the timing wedge in block $b$ is $w_{f,z,b,t}(h) = G^M_{f,z,b,t} -
R^{M,G}_{f,z,b,t}$, and
\begin{equation}
\label{eq:emissions}
E_{f,t} \;=\; \sum_{z \in \mathcal{Z}_f} \sum_{b \in \mathcal{B}}
\omega_b \cdot \mathbb{E}\big[\text{MCI}_{z,b,t}\big] \cdot w_{f,z,b,t}(h)
\;+\; \sum_{z} e^B_{f,z,t} \cdot B^{\text{run}}_{f,z,t},
\end{equation}
where $\omega_b$ is the share of period-$t$ hours in block $b$,
$\text{MCI}_{z,b,t}$ is the marginal carbon intensity of zone $z$ in block
$b$ (taken from EPA CAMD data and dispatch reconstruction), and the second
term is direct emissions from on-site backup running ($e^B$ is the emissions
rate; diesel has $e^B > 0$, storage has $e^B = 0$ if charged by renewables).

\paragraph{Hyperscaler Bellman equation.} The Markov-perfect value function
of firm $f$ satisfies
\begin{equation}
\label{eq:hyper_bellman}
W_f(\mathbf{s}_t) \;=\; \max_{\{a_{f,z,t}\}_{z \in \mathcal{Z}_f}} \mathbb{E}\left[
\pi_{f,t} \;+\; \beta\, W_f(\mathbf{s}_{t+1})
\,\Big|\, \sigma_{-f}, \mathbf{s}_t \right],
\end{equation}
where $\sigma_{-f}$ is the rivals' Markov strategy profile and the maximization
is over the action space (\ref{eq:hyper_action}) subject to feasibility
constraints stated in \S\ref{sec:feasibility} below.

\paragraph{Conditional separability across locations.} The hyperscaler payoff
is not separable across zones in general---the contest reward and the emissions
penalty both depend on choices at \emph{all} zones. It becomes separable once
we condition on two scalar aggregates and on prices. Writing $\lambda^H_{z,t}$
for the shadow price of interconnection headroom at $z$, $\lambda^Z_{z,t}$ for
the shadow price of the local clean-MW pool, and $p^{\text{aux}}_{z,t}$ for the
ancillary services price (all defined in \S\ref{sec:shortrun}), the per-period
payoff decomposes as
\begin{equation}
\label{eq:separability}
\pi_{f,t} \;=\; v_{f,t}(C_{f,t}) \;-\;
\sum_{z \in \mathcal{Z}_f} \tilde c_{f,z,t}\big(a_{f,z,t}; \lambda^H_z, \lambda^Z_z, p^{\text{aux}}_z\big)
\;-\; \Gamma\big(\gamma_f, E_{f,t}\big) V_t .
\end{equation}
The cost term in (\ref{eq:hyper_cost}) is additive across zones, so it poses no
obstacle. The two non-separable objects each enter through a scalar: the
contest reward $v_{f,t}$ depends on zone choices only through total compute
$C_{f,t} = \sum_{z}\sum_\sigma \kappa_z K_{f,z,t}^\sigma$
(eq.~\ref{eq:compute}), and the emissions penalty depends on them only through
total attributable emissions $E_{f,t}$ (eq.~\ref{eq:emissions}), itself a sum
across zones. Conditional on a candidate value of the pair
$(C_{f,t}, E_{f,t})$ and on the price/shadow-price vector, the per-zone
sub-problems decouple: each zone's contribution to $C$ and $E$ is fixed by its
own action, so each candidate location can be evaluated independently. This is
exactly the structure that the combinatorial discrete-choice method of
\citet{arkolakis2023combinatorial} exploits, and it lets us prune dominated
location sets without enumerating all $2^{|\mathcal{Z}_f|}$ combinations.

\paragraph{What the adaptation requires.} Two qualifications matter for
validity and are easy to elide. First, $v_{f,t}(C)$ is strictly concave in $C$
(eq.~\ref{eq:quality} with $\alpha_f<1$, eq.~\ref{eq:contest}), so the marginal
contest reward of compute is diminishing; the monotone-comparative-statics
argument that justifies pruning requires checking the single-crossing condition
in $C$ rather than assuming pure submodularity. Second, the resource
constraints (\ref{eq:headroom})--(\ref{eq:cleanpool}) couple \emph{firms}, not
just a firm's own zones, through $(\lambda^H_z,\lambda^Z_z)$. Separability is
therefore conditional on equilibrium shadow prices, and the location-pruning
step is nested inside the fixed-point iteration that clears those prices: we
prune given $(\lambda^H,\lambda^Z)$, re-clear the constraints given the implied
demands, and iterate to convergence. We do not adopt the pruning algorithm
``off the shelf''; we adopt it as the inner loop of a price-clearing
fixed point, and we verify the single-crossing condition numerically at the
estimated parameters. The formal statement and the proof that
(\ref{eq:separability}) holds under the maintained assumptions are deferred to
Appendix~B.

\paragraph{Feasibility.}
\label{sec:feasibility}
At each $(z,t)$, two resource constraints couple
firms' choices:
\begin{align}
\label{eq:headroom}
&\textbf{(Interconnection headroom)} & \sum_f \sum_\sigma K_{f,z,t}^\sigma
\;\leq\; H_{z,t}, \\
\label{eq:cleanpool}
&\textbf{(Clean-MW pool)} & \sum_f \big(K_{f,z,t}^{\text{PPA}} + K_{f,z,t}^{\text{Colo}}\big)
\;\leq\; Z_{z,t}.
\end{align}
Multipliers $\lambda^H_{z,t} \geq 0$ and $\lambda^Z_{z,t} \geq 0$ on these
constraints are the resource prices that enter the separable payoff
(\ref{eq:separability}). REC procurement does not draw against
(\ref{eq:cleanpool}) in the baseline because RECs trade in a regional
market; the counterfactual ``geographic REC restrictions'' makes
(\ref{eq:cleanpool}) bind for REC procurement as well.

\subsubsection{Generator Problem}
\label{sec:gen_lr}

Generation is owned by a set of firms $\mathcal{K} = \{1, \dots, K\}$
indexed by $k$. Each firm operates a portfolio of plants identified by fuel
type $l \in \mathcal{L}^{\text{fuel}}$ and location $z$. In each long-run
period, generator $k$ chooses for each $(l,z)$ pair an increment $\Delta
O^{\max}_{l,k,z,t} \in \mathbb{R}$ to installed capacity (positive = build,
negative = retire). Building incurs sunk cost $I^{\text{build}}_{l,z,t}$
financed at rate $1 + r_{l,k}(\sigma^2_{l,k})$, where the financing-rate
function $r_l(\cdot)$ is strictly increasing in net-revenue variance (a
reduced-form cost-of-capital channel). Retirement incurs $I^{\text{retire}}_{l,z,t}$
and recovers a salvage value $s^{\text{salv}}_{l,z,t}$.

\paragraph{Revenue streams.} A generator of type $l$ at zone $z$ earns
revenue from up to four streams: (i) energy revenue, $\bar p^{\text{ene}}_{l,z,t}
\cdot Q^{\text{ene}}_{l,k,z,t}$, where $\bar p^{\text{ene}}_{l,z,t}$ is the
LMP at zone $z$ and $Q^{\text{ene}}_{l,k,z,t}$ is the dispatched quantity
(both exogenous from the short-run market); (ii) ancillary-services revenue,
$\bar p^{\text{aux}}_{l,z,t} \cdot Q^{\text{aux}}_{l,k,z,t}$; (iii) capacity-market
revenue, $\bar p^{\text{cap}}_{z,t} \cdot O^{\max,\text{cleared}}_{l,k,z,t}$,
where $\bar p^{\text{cap}}$ is the cleared capacity-market price (from PJM
BRA and analogous mechanisms); and (iv) REC revenue for $l \in
\mathcal{L}^{\text{ren}}$, $p^{REC}_t \cdot \text{REC}^{\text{mint}}_{l,k,z,t}$.

\paragraph{PPA contracting.} Generators of renewable types may sign PPAs
with hyperscalers. A PPA between generator $k$ and hyperscaler $f$ at
$(z,t)$ specifies a fixed price $\bar p^{\text{PPA}}_{f,k,z}$ for a
quantity $\bar q^{\text{PPA}}_{f,k,z}$ over a tenor $T^{\text{PPA}}_{f,k}$.
Following the design specification, the PPA price is determined by Nash
bargaining over the joint surplus of $(f,k)$ relative to outside options:
\begin{equation}
\label{eq:ppa_bargain}
\bar p^{\text{PPA}}_{f,k,z} \;=\; \arg\max_{p}
\big[ U_f(p; \cdot) - U_f^{\text{out}} \big]^{\zeta_f}
\big[ \Pi_k(p; \cdot) - \Pi_k^{\text{out}} \big]^{1 - \zeta_f},
\end{equation}
where $\zeta_f \in [0,1]$ is hyperscaler $f$'s bargaining weight (a
structural parameter to be estimated from observed PPA prices),
$U_f^{\text{out}}$ is $f$'s outside option (next-best procurement mode at
$z$), and $\Pi_k^{\text{out}}$ is $k$'s outside option (merchant operation,
or PPA with a non-hyperscaler offtaker). The bargaining is over the
\emph{surplus from generator entry that is enabled by the contract}:
without the PPA, $k$ may not enter at all; with the PPA, $k$ enters and the
total surplus is divided.

\paragraph{Generator value function.} The generator's value function is
\begin{equation}
\label{eq:gen_bellman}
V_{k}(\mathbf{s}_t) \;=\; \max_{\{\Delta O^{\max}_{l,k,z,t}\}_{l,z}}
\mathbb{E}\left[ \pi^{\text{gen}}_{k,t} + \beta\, V_k(\mathbf{s}_{t+1})
\,\big|\, \mathbf{s}_t \right],
\end{equation}
with per-period payoff
\begin{equation}
\label{eq:gen_flow}
\pi^{\text{gen}}_{k,t} \;=\; \sum_{l,z}\Big[
\bar p^{\text{ene}}_{l,z,t} Q^{\text{ene}}_{l,k,z,t}
+ \bar p^{\text{aux}}_{l,z,t} Q^{\text{aux}}_{l,k,z,t}
+ \bar p^{\text{cap}}_{z,t} O^{\max,\text{cleared}}_{l,k,z,t}
+ p^{REC}_t \text{REC}^{\text{mint}}_{l,k,z,t}
\Big]
- C^{\text{gen}}_{k,t} - I^{\text{net}}_{k,t},
\end{equation}
where $C^{\text{gen}}_{k,t}$ is operating cost (variable cost times dispatched
quantity) and $I^{\text{net}}_{k,t}$ is net investment cost minus PPA
revenues over the period:
\begin{equation}
I^{\text{net}}_{k,t} = \sum_{l,z}\big[ I^{\text{build}}_{l,z,t} \mathbf{1}\{\Delta O^{\max}_{l,k,z,t} > 0\} \Delta O^{\max}_{l,k,z,t} (1 + r_l(\sigma^2_{l,k}))
- s^{\text{salv}}_{l,z,t} \mathbf{1}\{\Delta O^{\max}_{l,k,z,t} < 0\} |\Delta O^{\max}_{l,k,z,t}|\big]
- \sum_f \bar p^{\text{PPA}}_{f,k,z} \bar q^{\text{PPA}}_{f,k,z}.
\end{equation}

\paragraph{Free-entry condition for the marginal entrant.} For each
$(l, z, t)$ where new entry is observed, the marginal entrant earns
zero net present value. With $\Pi^{l,\sigma}_{k,z,t}$ denoting expected
discounted profit under contract structure $\sigma \in \{\text{merch, PPA,
Colo}\}$ for fuel $l$ at $(z,t)$,
\begin{equation}
\label{eq:gen_free_entry}
\Pi^{l, \sigma^*}_{k,z,t} \;\leq\; 0 \text{ for all } (l,k,z),\qquad
\Pi^{l, \sigma^*}_{k,z,t} \;=\; 0 \text{ if } \Delta O^{\max}_{l,k,z,t} > 0,
\end{equation}
where $\sigma^*$ is the equilibrium-selected contract structure. Because
$r_l$ is increasing in net-revenue variance and colocation eliminates that
variance, the standard ordering holds:
\begin{equation}
\label{eq:entry_ordering}
r_l(0) \;\leq\; r_l(\sigma^2_{\text{PPA}}) \;\leq\; r_l(\sigma^2_{\text{merch}})
\quad\Rightarrow\quad
\text{entry}_{\text{Colo}} \;\geq\; \text{entry}_{\text{PPA}} \;\geq\; \text{entry}_{\text{merch}}
\quad \text{(all else equal)}.
\end{equation}
This investment-hierarchy result is the central mechanism by which
hyperscaler procurement choices propagate into the supply stack: contracts
that reduce generator revenue variance lower the cost of capital and induce
entry. The empirical content of $r_l(\cdot)$ is recovered from observed
financing terms on contract-level data from S\&P Capital IQ Pro.


\section{Short-run}
\label{sec:shortrun}

The short-run is the within-period operations problem at $(z, t)$, conditional
on the long-run state $\mathbf{s}_t$. We partition each long-run period $t$
into a finite set of \emph{load blocks} $\mathcal{B} = \{1, \dots, B\}$ with
weights $\omega_b > 0$, $\sum_b \omega_b = 1$. Each block represents a class
of hours within $t$ characterized by similar load, renewable availability,
and ancillary services requirements (for instance, $B = 8$ blocks based on
$\{$peak/off-peak$\} \times \{$summer/winter/shoulder$\} \times $ renewable
availability tercile). The block structure simultaneously preserves the
timing wedge mechanism that drives REC compliance dynamics and keeps the
state space tractable.

Within each block, the short-run consists of four interlocking markets:
the energy market (\S\ref{sec:energy_market}), the REC market
(\S\ref{sec:rec_market}), the capacity market (\S\ref{sec:cap_market}),
and the ancillary services market (\S\ref{sec:aux_market}). The first
three are largely standard; the fourth is the central new construct of
the paper and the locus of identification via experimental data. The
hyperscaler demand side enters all four markets, and procurement-mode
costs (\S\ref{sec:proc_costs}) are derived as block-weighted expectations
of clearing prices.

\subsection{Demand Side}
\label{sec:demand}

Each hyperscaler $f$ at zone $z$ in block $b$ has gross IT load
\begin{equation}
\label{eq:gross_load}
D^M_{f,z,b,t} \;=\; \rho^{\text{train}}_{f,z,b,t} \cdot K^{\text{train}}_{f,z,t}
\;+\; \rho^{\text{infer}}_{f,z,b,t} \cdot K^{\text{infer}}_{f,z,t},
\end{equation}
where $K^{\text{train}}_{f,z,t}$ and $K^{\text{infer}}_{f,z,t}$ partition
$K_{f,z,t} = \sum_\sigma K_{f,z,t}^\sigma$ into training-allocated and
inference-allocated capacity, $\rho^{\text{train}}_{f,z,b,t}$ is the
utilization rate of training capacity in block $b$ (close to 1 for sustained
training), and $\rho^{\text{infer}}_{f,z,b,t}$ is the inference utilization
rate (block-dependent, reflecting diurnal user demand). The split between
training and inference at $z$ is a hyperscaler choice in the long-run problem;
in the baseline it is taken as observed (Aterio + Epoch + announced workload
mix), and the counterfactual ``relocate inference and training'' permutes
this split across zones.

\paragraph{Net grid draw.} Net grid draw by hyperscaler $f$ at $(z,b,t)$ is
\begin{equation}
\label{eq:grid_draw}
G^M_{f,z,b,t} \;=\; \big( D^M_{f,z,b,t}
- Q^{\text{contract}}_{f,z,b,t}
- Q^{\text{colo}}_{f,z,b,t}
- Q^{\text{stor,out}}_{f,z,b,t}
+ Q^{\text{stor,in}}_{f,z,b,t}
- B^{\text{run}}_{f,z,b,t} \big)_+,
\end{equation}
where $Q^{\text{contract}}$ is firm delivery under PPA, $Q^{\text{colo}}$ is
colocated generation, $Q^{\text{stor,out}}$ and $Q^{\text{stor,in}}$ are
storage discharge and charge respectively, and $B^{\text{run}}$ is backup
running. Storage dispatch obeys a standard energy balance constraint with
round-trip efficiency $\eta_S$:
\begin{equation}
\label{eq:storage_balance}
S^{\text{stock}}_{f,z,b+1,t} = S^{\text{stock}}_{f,z,b,t} + \eta_S \cdot
Q^{\text{stor,in}}_{f,z,b,t} - Q^{\text{stor,out}}_{f,z,b,t}, \quad
0 \leq S^{\text{stock}}_{f,z,b,t} \leq S_{f,z,t},
\end{equation}
and ramp constraints $|Q^{\text{stor,in/out}}_{f,z,b,t}| \leq \bar R_S \cdot
S_{f,z,t}$ where $\bar R_S$ is the (technology-determined) C-rate.

\paragraph{Load shape signatures.} Beyond gross MW, each block also has a
load-shape signature $\vec\zeta_{f,z,b,t}$ summarizing the dimensions of
the AI load that matter for grid impact beyond pure energy: ramp rate,
harmonic content (THD), power factor, transient response, and inertia
contribution. The vector $\vec\zeta$ is a function of the technical
parameters of the data center (chip generation, cooling architecture, power
electronics topology, batch scheduling, quantization), the load mix
(training vs.\ inference, model size), and the procurement/storage
configuration. The mapping from technical parameters to $\vec\zeta$ is
calibrated using the digital twin experiments described in
\S\ref{sec:experiments}; the role of $\vec\zeta$ in market clearing
appears in \S\ref{sec:aux_market}.

\subsection{RECs, PPAs, Colocation}
\label{sec:proc_costs}

This subsection derives the procurement-mode cost
$c^{\mathrm{ene}}_{z,\sigma,t}$ entering the long-run hyperscaler cost
(\ref{eq:hyper_cost}) from the short-run market clearing.

\paragraph{Merchant.} Under $\sigma = \text{merch}$ the hyperscaler buys
energy at LMP and retires no RECs:
\begin{equation}
\label{eq:c_merch}
c^{\mathrm{ene}}_{z,\text{merch},t} \;=\; \sum_{b \in \mathcal{B}} \omega_b
\, \mathbb{E}\big[ p^{\text{ene}}_{z,b,t} \big] \;+\; \lambda^H_{z,t},
\end{equation}
where $p^{\text{ene}}_{z,b,t}$ is the LMP in block $b$. The headroom shadow
price $\lambda^H_{z,t}$ enters because every mode draws against
(\ref{eq:headroom}).

\paragraph{REC.} Under $\sigma = \text{REC}$ the hyperscaler buys energy at
LMP and retires RECs sufficient to cover share $\phi$ of consumption subject
to granularity $h$. The granular and annual REC obligations are
\begin{align}
\label{eq:gran_req}
R^{M,G}_{f,z,b,t} &\geq h \phi D^M_{f,z,b,t}, \qquad
R^{M,G}_{f,z,b,t} \leq S^{\text{qual}}_{z,b,t}, \\
\label{eq:annual_req}
R^{M,A}_{f,z,t} &\geq (1 - h) \phi \sum_b \omega_b D^M_{f,z,b,t},
\end{align}
where $S^{\text{qual}}_{z,b,t} \leq S_{z,b,t}$ are RECs eligible by
geography and tier under prevailing rules. The per-MWh cost is
\begin{equation}
\label{eq:c_rec}
c^{\mathrm{ene}}_{z,\text{REC},t} \;=\; \sum_b \omega_b \mathbb{E}[p^{\text{ene}}_{z,b,t}]
\;+\; \phi\, \tilde p^{REC}_t(h) \;+\; \lambda^H_{z,t},
\end{equation}
where the effective shadow REC price $\tilde p^{REC}_t(h) = p^{REC}_t + \nu_t
+ \mu_t$ comprises the market REC price $p^{REC}_t$, the granular-constraint
multiplier $\nu_t$, and the annual-constraint multiplier $\mu_t$ (the latter
two derived in \S\ref{sec:rec_market}).

\paragraph{PPA.} Under $\sigma = \text{PPA}$ the hyperscaler receives
contracted quantity $\bar q^{\text{PPA}}_{f,k,z}$ at fixed price
$\bar p^{\text{PPA}}_{f,k,z}$ (set by Nash bargaining,
eq.~\ref{eq:ppa_bargain}) and meets residual demand at LMP:
\begin{equation}
\label{eq:c_ppa}
c^{\mathrm{ene}}_{z,\text{PPA},t} \;=\;
\frac{\bar q^{\text{PPA}}}{D^M_{f,z,t}} \bar p^{\text{PPA}}_{f,k,z}
+ \Big(1 - \frac{\bar q^{\text{PPA}}}{D^M_{f,z,t}}\Big) \sum_b \omega_b
\mathbb{E}[p^{\text{ene}}_{z,b,t}] + \lambda^H_{z,t} + \lambda^Z_{z,t},
\end{equation}
where $D^M_{f,z,t} = \sum_b \omega_b D^M_{f,z,b,t}$ is annual load. The
clean-MW shadow price $\lambda^Z_{z,t}$ enters because PPA draws against
(\ref{eq:cleanpool}). Under PPA the hyperscaler also receives the RECs
generated by the contracted output and retires them against
(\ref{eq:gran_req})--(\ref{eq:annual_req}).

\paragraph{Colocation.} Under $\sigma = \text{Colo}$ a dedicated generator
delivers $Q^{\text{colo}}_{f,z,b,t}$ behind the meter at a transfer price
$\tilde p^{\text{colo}}_{z,t}$. The hyperscaler's energy cost is
\begin{equation}
\label{eq:c_colo}
c^{\mathrm{ene}}_{z,\text{Colo},t} \;=\;
\frac{Q^{\text{colo}}}{D^M_{f,z,t}} \tilde p^{\text{colo}}_{z,t}
+ \Big(1 - \frac{Q^{\text{colo}}}{D^M_{f,z,t}}\Big) \sum_b \omega_b
\mathbb{E}[p^{\text{ene}}_{z,b,t}] + \lambda^H_{z,t} + \lambda^Z_{z,t}.
\end{equation}
Under colocation with storage, $Q^{\text{colo}}$ can be re-timed across
blocks subject to (\ref{eq:storage_balance}); this is the mechanism by
which colocation-with-storage substantively shrinks the timing wedge in
(\ref{eq:emissions}).

\subsection{REC Market Clearing}
\label{sec:rec_market}

Renewables mint one REC per MWh of dispatched generation in the block of
production: $r_{l,k,z,b,t} = \mathbf{1}\{l \in \mathcal{L}^{\text{ren}}\}
\cdot Q^{\text{ene}}_{l,k,z,b,t}$, and total block-level REC supply at
zone $z$ is $S_{z,b,t} = \sum_{l,k} r_{l,k,z,b,t}$. RECs may be banked
across blocks within $t$ subject to a stock cap $\bar K^{REC}$ and a
banking depreciation $\delta^{REC}$:
\begin{equation}
\label{eq:rec_bank}
K^{REC}_{z,b+1,t} \;=\; (1 - \delta^{REC}) K^{REC}_{z,b,t}
+ S_{z,b,t} - R^{\text{ret}}_{z,b,t},
\qquad K^{REC}_{z,b,t} \in [0, \bar K^{REC}].
\end{equation}
We suppress cross-period (across $t$) REC banking to preserve separability
between the long-run capacity decisions and the short-run market clearing;
this is a tractability simplification that suppresses an empirically
small banking margin observed in practice.

\paragraph{REC market clearing problem.} The competitive REC allocation
minimizes the present cost of compliance given hyperscaler obligations
(\ref{eq:gran_req})--(\ref{eq:annual_req}) and the REC banking constraint
(\ref{eq:rec_bank}):
\begin{equation}
\label{eq:rec_problem}
\min_{\{R^{\text{ret}}, K^{REC}, R^{M,G}, R^{M,A}\}}
\sum_b \omega_b \big[ p^{REC}_t R^{\text{ret}}_{z,b,t}
+ \bar p^{\text{ACP}}_t \cdot \text{short}_{z,b,t} \big],
\end{equation}
subject to (\ref{eq:gran_req}), (\ref{eq:annual_req}), (\ref{eq:rec_bank}),
the supply-availability constraint
$R^{\text{ret}}_{z,b,t} \leq S^{\text{qual}}_{z,b,t} + K^{REC}_{z,b,t}$,
and the non-negativity of $R^{\text{ret}}_{z,b,t}$ and shortfalls
$\text{short}_{z,b,t} \geq 0$, where $\bar p^{\text{ACP}}_t$ is the
alternative compliance payment (state-RPS-specific ceiling). Letting
$\nu_{z,b,t}$ be the multiplier on the granular constraint
(\ref{eq:gran_req}) and $\mu_{z,t}$ the multiplier on the annual
constraint (\ref{eq:annual_req}), the effective shadow REC price entering
(\ref{eq:c_rec}) is
\begin{equation}
\label{eq:rec_shadow}
\tilde p^{REC}_t(h) \;=\; p^{REC}_t + \nu_{z,b,t} + \mu_{z,t},
\end{equation}
with $p^{REC}_t \leq \bar p^{\text{ACP}}_t$, equality when shortfalls
occur. Interior banking yields the standard Euler condition
$p^{REC}_t = \beta(1 - \delta^{REC}) \mathbb{E}_t[p^{REC}_{t+1}]$ for
the cross-period market.

\paragraph{REC market clearing is regional.} Unlike the resource
constraints (\ref{eq:headroom})--(\ref{eq:cleanpool}), the REC market
clears at the regional level: $\sum_b \omega_b \sum_z S^{\text{qual}}_{z,b,t}
\geq \sum_b \omega_b \sum_z \sum_f R^{M,G}_{f,z,b,t} + \sum_z \sum_f
R^{M,A}_{f,z,t}$, with the regional aggregation following the prevailing
geographic-qualification rule. The counterfactual ``geographic REC
restrictions'' tightens the qualification rule at the zone level; the
counterfactual ``national REC market'' loosens it to a single national
pool.

\subsection{Energy Market}
\label{sec:energy_market}

The energy market is solved \emph{exogenously}. For each $(z,b,t)$ the LMP
$p^{\text{ene}}_{z,b,t}$ is taken from observed ISO data (PJM Day-Ahead,
ERCOT, CAISO, MISO, SPP, ISO-NE, NYISO) for in-sample periods and from
ISO-published cost curves combined with futures prices for the simulation
window. Formally, write $\mathcal{C}^{\text{stack}}_{z,b,t}(\cdot)$ for the
zone's marginal cost stack (a step function of dispatched quantity) and
$D^{\text{tot}}_{z,b,t}$ for total load (hyperscaler plus residual). The
LMP is
\begin{equation}
\label{eq:lmp}
p^{\text{ene}}_{z,b,t} \;=\; \mathcal{C}^{\text{stack}}_{z,b,t}\big(D^{\text{tot}}_{z,b,t}\big),
\end{equation}
with transmission-constrained LMP differences across zones taken from
historical data. In counterfactuals, the cost stack is shifted to reflect
endogenous generator-capacity choices: if a counterfactual generates
$\Delta O^{\max}_{l,k,z,t}$ MW of new fuel-$l$ entry at zone $z$, the cost
stack is updated by inserting (or removing) a block of capacity at the
fuel's variable cost.

\paragraph{Why exogenous.} The rationale and its cost are stated once, for all
imported price processes, in the endogeneity/exogeneity overview at the head of
\S\ref{sec:model}: endogenizing the LMP would require a unit-commitment plus
optimal-power-flow solve at each estimation iteration, which is infeasible at
this scope. The point specific to the energy market is one of magnitude. The
LMP feedback \emph{within} a period---hyperscaler siting nudging the same
period's clearing price---is second-order relative to the LMP feedback
\emph{across} years that operates through generator entry and exit, and it is
the latter that drives the long-run counterfactuals. Treating the within-period
loop as open therefore discards little of the variation the counterfactuals
turn on.

\subsection{Capacity Market}
\label{sec:cap_market}

For zones within ISOs that operate centralized capacity markets (most
prominently PJM via the Base Residual Auction, but also ISO-NE Forward
Capacity Market and NYISO ICAP), the capacity-market clearing price
$\bar p^{\text{cap}}_{z,t}$ is taken exogenously from auction results.
The capacity-market revenue stream affects generator entry through the
generator value function (\ref{eq:gen_bellman}). Two features of PJM's
capacity construct are central to the AI-grid question:

\paragraph{Capacity Performance.} PJM's Capacity Performance (CP) construct
requires capacity-market participants to deliver during emergencies, with
substantial penalties for nonperformance. Data centers with on-site
generation or storage can in principle qualify as demand-response or
behind-the-meter capacity resources, generating capacity-market revenue
that offsets infrastructure cost. We model this via an additional revenue
term for the hyperscaler:
\begin{equation}
\label{eq:hyper_cap_rev}
\text{CapRev}_{f,z,t} \;=\; \bar p^{\text{cap}}_{z,t} \cdot \min\big\{
S_{f,z,t} \cdot \xi^{\text{cap}}_{S} \;,\; \overline{\text{DR}}_{f,z,t}
\big\},
\end{equation}
where $\xi^{\text{cap}}_S$ is the storage-to-capacity derating factor and
$\overline{\text{DR}}_{f,z,t}$ is the demand-response cap (technology- and
zone-dependent). The mechanism is that storage installed for outage
protection can simultaneously bid into the capacity market, lowering the
effective cost of storage; this is a counterfactual lever of interest
because it interacts with the BYOC mandate scenario.

\paragraph{Capacity accreditation.} A growing share of capacity markets are
moving to effective-load-carrying-capability (ELCC) accreditation, which
discounts intermittent resources. We let the cleared capacity
$O^{\max,\text{cleared}}_{l,k,z,t}$ in the generator revenue equation
(\ref{eq:gen_flow}) be the ELCC-derated value, with derating factors taken
from ISO publications and held fixed in the baseline.

\subsection{Ancillary Services Market}
\label{sec:aux_market}

This is the central new market in the model. Ancillary services---regulation
up/down, spinning reserve, non-spinning reserve, supplemental reserve,
voltage support/reactive power, frequency response, and emerging products
for inertia and fast-frequency response---are the segment of wholesale
electricity in which the distinctive technical features of AI load have the
largest disproportionate impact, both because AI ramps are fast and large
relative to traditional industrial load and because AI power electronics
inject harmonics and demand reactive power in patterns that differ from
conventional loads.

\paragraph{Why ancillary services are different.} Three features distinguish
the ancillary services market from energy and capacity:
\begin{enumerate}[leftmargin=*]
    \item \emph{Co-optimization.} Most U.S.\ ISOs co-optimize energy and
    ancillary services in the day-ahead and real-time markets, so the same
    generator can provide both. Ancillary-services prices reflect the
    opportunity cost of foregone energy revenue plus the technical cost of
    providing the service.
    \item \emph{Heterogeneous products.} Ancillary services are not a
    single good. PJM operates Regulation, Synchronized Reserve, Primary
    Reserve, and 30-Minute Reserve, each with distinct technical
    requirements. ERCOT operates Regulation Up/Down, Responsive Reserve,
    Non-Spinning Reserve, and ERCOT Contingency Reserve. We index ancillary
    services products by $j \in \mathcal{J}^{\text{aux}}$.
    \item \emph{Demand is determined by load characteristics, not just
    magnitude.} The ISO sets ancillary-services requirements partly as a
    function of load forecast uncertainty and the load's contribution to
    ramping needs. Two MW of constant industrial load and two MW of bursty
    AI inference load generate different ancillary-services requirements.
    This is the central econometric question of the paper.
\end{enumerate}

\paragraph{Aggregate demand for ancillary services.} For each product $j$
and zone $z$ in block $b$, total ancillary-services demand is
\begin{equation}
\label{eq:aux_demand}
D^{\text{aux}}_{j,z,b,t} \;=\;
\underbrace{\bar D^{\text{aux},0}_{j,z,b,t}}_{\text{baseline (ISO rule)}}
\;+\; \sum_f
\underbrace{\Delta^{\text{aux}}_{j}\big(\vec\zeta_{f,z,b,t}, K_{f,z,t}\big)}_{\text{AI contribution}}
\;+\; \xi^{\text{aux}}_{j,z,b,t},
\end{equation}
where $\bar D^{\text{aux},0}$ is the baseline ancillary-services requirement
under the ISO's rule (taken from ISO documentation), $\Delta^{\text{aux}}_j(\cdot)$
is the marginal ancillary-services demand contributed by hyperscaler $f$ at
$(z,b,t)$, $\vec\zeta_{f,z,b,t}$ is the load-shape signature from
\S\ref{sec:demand}, and $\xi^{\text{aux}}_{j,z,b,t}$ is an exogenous
non-AI residual.

The function $\Delta^{\text{aux}}_j$ is the central econometric object of
the paper. It maps technical features of the AI load (captured in
$\vec\zeta$) into incremental ancillary-services demand. We do not impose
this function from theory; we estimate it.

\paragraph{Identification of $\Delta^{\text{aux}}_j$.} We pursue two
identification strategies that exploit different sources of variation and
serve as cross-checks on each other.

\emph{Strategy 1: Observational DiD on AI model releases.} Following our
prior reduced-form work, we use the staggered timing of AI model release
dates as the treatment shock. The treatment exposure of zone $z$ is built
from the share of zone-$z$ data centers operated by the releasing firm
(measured via Aterio) and the inferred training/inference allocation
(Epoch + announcements). The outcome is the zone-block-product ancillary
services price $p^{\text{aux}}_{j,z,b,t}$. A stacked DiD recovers the
average treatment effect of an AI model release on ancillary-services
prices, which combined with the (exogenous, ISO-published) ancillary-services
supply curve yields the implied shift in ancillary-services demand.
The estimating equation is
\begin{equation}
\label{eq:aux_did}
\log p^{\text{aux}}_{j,z,b,t,m} \;=\; \sum_{k=-K}^K \beta^{j}_k D^k_{z,t,m}
+ \gamma_{z,m} + \delta_{t,m} + \mathbf{X}'_{z,t}\theta + \varepsilon_{z,t,m},
\end{equation}
where $m$ indexes model-release events and $D^k_{z,t,m}$ is the standard
stacked-DiD event-time treatment indicator.

\emph{What identifies $\beta^j$.} The design is a shift-share: the ``shock'' is
a national, firm-level event (a model release by firm $f$) whose timing is
plausibly unrelated to any single zone's ancillary-services market, and the
``share'' is zone $z$'s differential exposure to that firm, built from its
data-center mix and inferred training/inference allocation. Identification
rests on two restrictions. \emph{(i) Conditional quasi-randomness of exposure:}
within an event window, after $\gamma_{z,m}$ (zone-by-event) and $\delta_{t,m}$
(time-by-event) fixed effects and controls $\mathbf{X}$ for weather, fuel
prices, and residual (non-AI) load, the more-exposed and less-exposed zones
would have followed parallel ancillary-services-price paths absent the release.
\emph{(ii) Exclusion:} a model release affects $p^{\text{aux}}$ \emph{only}
through the AI load it induces in exposed zones, not through a contemporaneous
channel (e.g., a market-wide capital-expenditure cycle) correlated with release
timing. Restriction~(i) is testable through the pre-period event-time
coefficients and the placebo and alternative-estimator suite of
\S\ref{sec:est_short}; restriction~(ii) is the maintained assumption, and is
the reason the digital-twin strategy below is valuable as an independent
cross-check rather than a mere robustness exercise.

\emph{From a price effect to a demand shift.} Ancillary-services demand is set
administratively by the ISO requirement plus the AI contribution
(eq.~\ref{eq:aux_demand}); to first order it is price-inelastic over the
relevant range. A model release therefore acts as a horizontal demand shift
$\Delta D^{\text{aux}}$ that the market absorbs by moving up the
(upward-sloping) supply curve. With supply elasticity
$\varepsilon^s_{j,z,b,t} = (\partial S/\partial p)(\bar p/\bar Q)$ evaluated at
the pre-shock clearing point $(\bar p^{\text{aux}}_{j,z,b,t},
\bar Q^{\text{aux}}_{j,z,b,t})$, a quantity shift $\Delta D$ raises the log
price by $\Delta\log p \approx (1/\varepsilon^s)\,(\Delta D/\bar Q)$. Inverting
and identifying $\hat\beta^j_k \approx \Delta\log p^{\text{aux}}$:
\begin{equation}
\label{eq:aux_demand_from_price}
\widehat{\Delta D^{\text{aux}}_{j,z,b,t}} \;=\; \hat\beta^j_k \cdot
\varepsilon^s_{j,z,b,t} \cdot \bar Q^{\text{aux}}_{j,z,b,t},
\end{equation}
where $\bar Q^{\text{aux}}_{j,z,b,t}$ is the pre-shock \emph{cleared quantity}
of product $j$ (in MW), not the price---the implied demand shift carries units
of quantity, so the level on the right must be a quantity. The supply elasticity
$\varepsilon^s$ is estimated separately from the slope of the ISO's posted
ancillary-services supply curve at the relevant operating point.\footnote{If
ancillary-services demand is not perfectly inelastic but has its own price
semi-elasticity $\varepsilon^d \le 0$, the price move reflects the gap between
the two slopes and (\ref{eq:aux_demand_from_price}) generalizes to
$\widehat{\Delta D} = \hat\beta^j_k\,\bar Q\,(\varepsilon^s-\varepsilon^d)$. We
treat the inelastic case ($\varepsilon^d=0$) as the baseline because the ISO
requirement is a rule, not a bid, and report the general case as a robustness
check.}

\emph{Strategy 2: Digital twin experiments.} The observational strategy
identifies the average effect of an AI model release on a treated zone but
cannot decompose that effect into contributions from the technical
parameters of the load (model size, quantization, batch scheduling, storage
configuration). The digital twin experiments described in
\S\ref{sec:experiments} provide exactly this decomposition by varying
single technical parameters in a controlled environment. Specifically,
for each technical parameter $\theta^{\text{tech}}_\ell$ (e.g., quantization
level, batch size, power-factor-correction topology) varied in experimental
trial $\ell$, the digital twin produces measurements of $\vec\zeta$. We
estimate
\begin{equation}
\label{eq:zeta_map}
\vec\zeta = \Lambda(\theta^{\text{tech}}; X^{\text{exp}}) + \eta,
\end{equation}
where $X^{\text{exp}}$ is the experimental design state (microgrid
configuration, ambient conditions) and $\eta$ is measurement noise.

\emph{What the twin identifies, and what it does not.} The twin operates at
four GPUs; hyperscaler zones operate at hundreds of megawatts. The bridge
between them is a division of labor between the \emph{intensive} and
\emph{extensive} margins of the load. The objects $\vec\zeta$ are deliberately
\emph{per-MW, intensive-margin} signatures---ramp rate per MW, total harmonic
distortion, displacement power factor, transient response per unit step---that
are properties of the load's \emph{shape}, fixed by chip generation, converter
topology, control firmware, and workload schedule, and approximately invariant
to how many identical units are racked behind the same point of common
coupling. What the twin identifies is therefore the \emph{shape map}
$\Lambda$: how a marginal change in a technical parameter moves the per-MW
signature, holding scale fixed. What the twin does \emph{not} identify is the
extensive margin---total installed MW, the fleet mix of technical parameters,
and the cross-unit aggregation of harmonics (which can attenuate through
phase cancellation or amplify through resonance). The extensive margin is
supplied by observational data: total MW from Aterio, technical mix inferred
from Epoch and announcements, and the realized aggregate effect from the DiD.
The twin supplies the derivative; the field data supplies the level. The
policy-relevant impact $\Delta^{\text{aux}}_j$ is then the composition
\begin{equation}
\label{eq:aux_composition}
\Delta^{\text{aux}}_j\big(\vec\zeta_{f,z,b,t}, K_{f,z,t}\big)
\;=\; g_j\!\Big(\Lambda(\theta^{\text{tech}}_{f,z,t}; X^{\text{exp}})\Big)
\cdot K_{f,z,t} \cdot a_{j,z}\big(\textstyle\sum_{f'} K_{f',z,t}\big),
\end{equation}
where $g_j(\cdot)$ maps a per-MW signature into per-MW ancillary-services
demand for product $j$ (the engineering relationship, e.g., reserve scaled by
ramp-rate exceedance), $K_{f,z,t}$ is the extensive scale, and
$a_{j,z}(\cdot)\in(0,1]$ is an aggregation factor capturing cross-unit
cancellation/resonance at the zone, calibrated where possible against the
utility-level Whisker Labs THD distribution (\S\ref{sec:data}) and otherwise
set to unity with the resulting bias signed and bounded. Separating the three
factors is what makes a four-GPU experiment informative about a gigawatt zone:
only $g_j(\Lambda(\cdot))$ is asked of the twin, and it is the one factor that
is scale-free.

\emph{Cross-validation.} The two strategies meet at the zone level. Aggregating
the twin-implied demand---inserting the field-measured scale $K_{f,z,t}$ and
mix into (\ref{eq:aux_composition}) and summing over firms---produces a
predicted $\sum_f \Delta^{\text{aux}}_j$ that should reproduce the
DiD-implied $\widehat{\Delta D^{\text{aux}}_{j,z,b,t}}$ of
(\ref{eq:aux_demand_from_price}). The comparison is a genuine overidentifying
test, not a calibration: the twin never sees ancillary-services prices, and the
DiD never sees a technical parameter, so agreement is informative.
Discrepancies are diagnostic and signed. If the twin under-predicts, either it
omits a feature of full-scale operation (cross-unit resonance, transformer and
cabling effects absent at four GPUs) or the aggregation factor $a_{j,z}$ is
mis-set; if it over-predicts, the observational variation likely reflects
factors beyond hyperscaler load and the DiD exclusion restriction needs
reconsideration. We adopt the convention that when the two agree within their
confidence intervals we use the twin-based composition for counterfactuals
(because only it provides the within-load decomposition needed to evaluate
technical-parameter policies); when they disagree we report both, treat the DiD
as the binding aggregate, and reconcile the gap through $a_{j,z}$.

\paragraph{Supply of ancillary services.} The supply side of the ancillary
services market is solved exogenously using ISO-published cost curves and
historical clearing data. For each product $j$ at $(z,b,t)$ the ancillary
services price is
\begin{equation}
\label{eq:aux_price}
p^{\text{aux}}_{j,z,b,t} \;=\; \mathcal{C}^{\text{aux,stack}}_{j,z,b,t}\big(D^{\text{aux}}_{j,z,b,t}\big),
\end{equation}
where $\mathcal{C}^{\text{aux,stack}}$ is the supply stack constructed from
generator offer curves (or the ISO's posted clearing data when offer
detail is unavailable). The exogenous treatment of supply is consistent
with the energy-market treatment in \S\ref{sec:energy_market}; the same
caveat applies (within-period feedback from demand to supply is muted,
across-period feedback operates through entry/retirement).

\paragraph{Hyperscaler attribution of ancillary-services cost.} The hyperscaler
ancillary-services cost in (\ref{eq:hyper_cost}) is
\begin{equation}
\label{eq:hyper_aux_cost}
c^{\mathrm{aux}}_{z,t} \cdot u^{\mathrm{aux}}_{f,z,t} \;=\;
\sum_j \sum_b \omega_b \cdot p^{\text{aux}}_{j,z,b,t} \cdot
\frac{\Delta^{\text{aux}}_j(\vec\zeta_{f,z,b,t}, K_{f,z,t})}{D^{\text{aux}}_{j,z,b,t}}
\cdot D^{\text{aux}}_{j,z,b,t},
\end{equation}
which is the firm $f$'s pro-rata share of total ancillary-services costs at
$(z,b,t)$. The attribution rule matters for counterfactuals on
cost-causation: in the baseline we use proportional attribution (consistent
with most ISO cost-allocation rules), but the counterfactual ``ancillary
services market design'' can shift this to incremental attribution (charge
the marginal contributor at the marginal price).

\subsection{Storage as a Dual-Use Asset}
\label{sec:storage_dual}

The hyperscaler's storage choice $S_{f,z,t}$ enters the model through
three channels, which the dual-use formulation makes explicit:

\begin{enumerate}[leftmargin=*]
    \item \emph{Outage protection.} Storage reduces the outage probability
    $P^{\mathrm{out}}$ in (\ref{eq:outage_cost}), lowering expected outage
    cost.
    \item \emph{Timing-wedge reduction.} Storage allows re-timing of
    behind-the-meter renewable generation under colocation, shrinking the
    timing wedge $w_{f,z,b,t}(h)$ and reducing attributable emissions in
    (\ref{eq:emissions}).
    \item \emph{Ancillary services participation.} Battery storage can
    provide regulation and reserve services. When the hyperscaler enrolls
    storage in the ancillary services market, the
    $\Delta^{\text{aux}}_j(\vec\zeta, K, S)$ function decreases (storage
    contributes to supply rather than demand), and the hyperscaler earns
    revenue $\text{AuxRev}_{f,z,t} = \sum_j \sum_b \omega_b p^{\text{aux}}_{j,z,b,t}
    Q^{\text{aux,supply}}_{f,z,b,t}$.
\end{enumerate}

The three channels typically pull in the same direction (storage is
valuable), but their relative magnitudes depend on parameters that the
digital twin experiments are well-positioned to recover. In particular,
the ancillary-services-participation channel depends on storage's
technical capability to ramp at sub-second time scales, which is a
function of inverter design and control software---both observable in the
twin experiments but not in observational data.

\subsection{Timing and Equilibrium}
\label{sec:timing}

The timing within a long-run period $t$ is:
\begin{enumerate}[leftmargin=*,label=\textbf{(\alph*)}]
    \item \emph{Long-run actions.} Hyperscalers and generators simultaneously
    choose long-run actions: hyperscaler $(\Delta K, \sigma, \Delta S, B)$
    via (\ref{eq:hyper_action}), generator $\Delta O^{\max}$. PPA terms are
    set by Nash bargaining (\ref{eq:ppa_bargain}). Resource constraints
    (\ref{eq:headroom})--(\ref{eq:cleanpool}) clear via shadow prices.
    \item \emph{Short-run market clearing.} Within each block $b$, the
    energy, capacity, ancillary services, and REC markets clear taking
    long-run capacities as given.
    \item \emph{Settlement.} Hyperscaler payoffs $\pi_{f,t}$ and generator
    payoffs $\pi^{\text{gen}}_{k,t}$ are realized.
\end{enumerate}

The equilibrium concept is a Markov-perfect equilibrium of the long-run
game in which (i) each hyperscaler's strategy attains the maximum in
(\ref{eq:hyper_bellman}) taking rivals' strategies and short-run prices as
given; (ii) each generator's strategy attains the maximum in
(\ref{eq:gen_bellman}); (iii) for each long-run period the short-run
markets clear; (iv) free entry (\ref{eq:gen_free_entry}) holds for generators;
(v) PPA prices satisfy the Nash-bargaining condition; (vi) investor
expectations in the contest (\ref{eq:contest}) are consistent with the
equilibrium quality distribution; and (vii) for each $(z,t)$ resource
prices $(\lambda^H_{z,t}, \lambda^Z_{z,t})$ clear the resource constraints.

Existence under sunk capacity and discrete actions follows from standard
arguments for dynamic discrete games with capacity as state
\citep{gowrisankaran2012capacity}. Uniqueness is not guaranteed; this is
consistent with the partial-identification posture we adopt in estimation.


\section{Data}
\label{sec:data}

The empirical implementation draws on a combination of administrative,
proprietary, and experimental data sources, organized below by the model
component each informs.

\subsection{Hyperscaler and Generator Data}

\paragraph{Data center locations and capacities.} The primary source for
the hyperscaler capacity stock $\{K_{f,z,t}^\sigma\}$ is the Aterio
data-center registry \citep{aterio_us_data_centers_2025}, which provides
facility-level location, operator, declared MW capacity, build status, and
announced expansions for the U.S.\ universe of data centers. We
cross-validate Aterio coverage against S\&P 451/DC~Bytes and against
state-level interconnection-queue filings.

\paragraph{Procurement contracts and financing.} PPA timing, counterparty,
tenor, contracted capacity, buyer type, and contract form (physical vs.\
virtual) are drawn from S\&P Capital IQ~Pro, supplemented by SEC filings and
the LevelTen PPA index for benchmark prices. Project- and firm-level asset
transactions and financing events around FID and COD are also from Capital
IQ~Pro. These data identify the contracting margin separately from the
entry margin: the same project can appear with FID before or after PPA
signing, which lets us isolate whether the PPA enables entry that would
not have occurred under merchant operation.

\paragraph{AI model releases and compute.} AI model release dates, training
compute estimates, and model release windows are from Epoch and from
public corporate announcements. These data construct the empirical analog
of model quality $q_{f,t}$ and provide the staggered treatment timing for
the DiD identification of ancillary-services impacts.

\paragraph{Generator data.} Plant-level capacity, fuel type, operator,
heat rate, and retirement status come from EIA Form~860 and EIA Form~923.
Day-ahead generation, fuel use, and emissions come from EPA CAMD
\citep{EPA_CAPMD_PowerSector2022} and from ISO-published unit-level
dispatch records where available.

\subsection{Market Data}

\paragraph{Energy market.} LMPs at the nodal and zonal level come from each
ISO's open data portal (PJM Data Miner, ERCOT MIS, CAISO OASIS, MISO,
NYISO, ISO-NE, SPP). For zones outside organized wholesale markets
(particularly the SERC and Western non-ISO footprints) we use EIA Form 930
hourly load and reconstructed marginal costs.

\paragraph{Capacity market.} PJM Base Residual Auction results, ISO-NE
Forward Capacity Auction results, and NYISO ICAP auction results provide
the cleared capacity-market prices $\bar p^{\text{cap}}_{z,t}$. ELCC
derating factors are from each ISO's accreditation rules.

\paragraph{Ancillary services market.} Ancillary-services clearing prices
and cleared quantities by product and zone come from each ISO's
ancillary-services market reports. PJM publishes hourly Synchronized
Reserve and Regulation clearing data; ERCOT publishes settlement-point
data for its four ancillary-services products; CAISO publishes
ancillary-services clearing by sub-region. These data define the supply
curve $\mathcal{C}^{\text{aux,stack}}$ entering (\ref{eq:aux_price}) and
serve as the dependent variable in the DiD identification
(\ref{eq:aux_did}).

\paragraph{REC market.} REC prices come from the S\&P Global Platts
Renewable Energy Price database, the M-RETS and PJM-GATS issuance and
retirement records, and state-level compliance market reporting.

\subsection{Network Flow Data}

EIA Form~930 provides hourly net generation, interchange, and load by
balancing authority, which we use to estimate the residual non-hyperscaler
demand $D^O_{z,t}$ and to construct the inter-echelon arc flows. FERC
Form~715 transmission-planning data establish the topology of the directed
graph and the transfer capacities $\bar T_{a,t}$. Where Form~715 is
aggregated at a higher spatial resolution than our zone definitions, we
allocate transfer capacity in proportion to the constituent zones' load.

\subsection{Power Quality Data}

Power-quality observations are from Whisker Labs' THD/voltage-sag
distribution at the utility-month level \citep{WhiskerLabs_THD_2024,
LaughnerEtAl_2024_THDonUSGrid}. These data validate the digital-twin
$\Lambda(\cdot)$ map at the aggregate level: predicted aggregate THD
contributions from data center load (constructed by aggregating
$\vec\zeta$ across hyperscaler load at each utility) should reproduce
the observed THD pattern.

\subsection{Experimental Data}

Experimental data from the CEWIT/AERTC digital twin (\S\ref{sec:experiments})
provide the $\Lambda(\cdot)$ calibration and serve as input to the
ancillary-services demand decomposition.

\section{Estimation Strategy}
\label{sec:estimation}

Estimation proceeds in three stages, exploiting the model's separation
between the long-run hyperscaler/generator problem, the short-run market
clearing, and the experimental calibration.

\subsection{Stage 1: Short-run Estimation}
\label{sec:est_short}

\paragraph{Ancillary-services DiD.} We estimate (\ref{eq:aux_did}) for each
ancillary-services product $j$ using model releases as the treatment shock
and the stacked-DiD design developed in our prior work to handle staggered
adoption. Treatment exposure is constructed at the zone-month level from
Aterio coverage of data centers and the inferred AI workload share at
each facility (using Epoch and announcements). Controls include weather,
fuel-price moves, residual (non-AI) load, and ISO-fixed-effect interactions
with month-of-sample to absorb market-wide ancillary-services trends. We
report pre-trend tests, placebo exercises with randomized treatment
assignment, the Callaway-Sant'Anna estimator as an alternative
identification strategy \citep{callaway2021difference}, and the
Borusyak-Jaravel-Spiess imputation estimator \citep{borusyak2024revisiting}
as robustness checks, mirroring the validation suite from our prior
reduced-form work.

\paragraph{Implied $\Delta^{\text{aux}}$.} The DiD coefficient
$\hat\beta^j_k$ on ancillary-services price is mapped to an
ancillary-services demand quantity via (\ref{eq:aux_demand_from_price}),
using the ISO-published ancillary-services supply curve slope at the
relevant clearing point. The implied $\widehat{\Delta D^{\text{aux}}_{j,z,b,t}}$
is then a moment that the experimental calibration must match in
aggregate.

\paragraph{Other short-run prices.} LMPs, capacity prices, and REC prices
are imported as data and are not separately estimated.

\subsection{Stage 2: Experimental Calibration}
\label{sec:est_exp}

The digital-twin experiments produce trial-level data $(X^{\text{exp}}_\ell,
\theta^{\text{tech}}_\ell, \vec\zeta_\ell, \eta_\ell)$. We estimate the
$\Lambda$ map in (\ref{eq:zeta_map}) by partially-linear regression of each
component of $\vec\zeta$ on $\theta^{\text{tech}}$, controlling
nonparametrically for the design state $X^{\text{exp}}$. The result is a
fitted function $\hat\Lambda$ that maps technical parameters to load-shape
signatures.

\paragraph{Composition with ancillary-services demand.} We then estimate
the partial $\Delta^{\text{aux}}_j(\vec\zeta, K)$ at the within-facility
level using $\hat\Lambda$ to assign $\vec\zeta$ to counterfactual
technical-parameter values, and we use the aggregated implied
ancillary-services demand as a moment to match the DiD-implied
$\widehat{\Delta D^{\text{aux}}}$ from Stage~1.

\paragraph{Cross-validation.} The aggregate match between
twin-implied and observationally-implied $\Delta^{\text{aux}}$ is the
diagnostic for the integration. We do not assume the twin is correct;
we test it.

\subsection{Stage 3: Long-run Estimation}
\label{sec:est_long}

The long-run hyperscaler problem is estimated by moment inequalities in
the spirit of \citet{holmes2011walmart}. The identifying restriction is
that observed firm portfolios dominate small unilateral deviations in
expected discounted profit:
\begin{equation}
\label{eq:moment_ineq}
\mathbb{E}\!\left[W_f(\mathbf{s}_t; a_{f,t}) - W_f(\mathbf{s}_t; a'_{f,t})
\,\Big|\, \mathbf{X}_t \right] \geq 0,
\end{equation}
for any feasible deviation $a'_{f,t}$ and information set $\mathbf{X}_t$.
We construct four types of deviations:
\begin{enumerate}[leftmargin=*]
    \item \emph{Location swaps.} Replace an opened zone $z$ with an
    unopened candidate $z'$ at the same capacity and procurement mode.
    Identifies relative zone attractiveness.
    \item \emph{Procurement swaps.} Hold $(z,k)$ fixed and vary $\sigma \in
    \mathcal{I}$. Identifies $\Gamma$, $r_l$, and the relative cost of
    procurement modes.
    \item \emph{Capacity swaps.} Adjust $k$ at a chosen zone. Identifies
    $\alpha_f$ and $\eta$ from the contest function.
    \item \emph{Storage swaps.} Adjust $\Delta S$ at a chosen zone.
    Identifies the marginal value of storage across its three channels
    (outage protection, timing-wedge reduction, ancillary-services
    participation).
\end{enumerate}
The deviation set is constructed to be feasible under
(\ref{eq:headroom})--(\ref{eq:cleanpool}) given the observed equilibrium
rival capacity. Standard errors and confidence sets follow the
partial-identification inference of \citet{andrews2010inference} and
related approaches.

\paragraph{Generator entry.} The generator free-entry condition
(\ref{eq:gen_free_entry}) provides moment conditions on the cost of
entry $I^{\text{build}}_{l,z,t}$ and the financing rate $r_l(\sigma^2)$.
We estimate these by matching observed entry rates by fuel type, zone,
and contract structure to model-predicted entry rates under the
estimated demand-side parameters.

\paragraph{PPA bargaining weight.} The bargaining weight $\zeta_f$ in
(\ref{eq:ppa_bargain}) is identified from observed PPA prices relative to
the model-implied outside options of each contracting pair.

\paragraph{Identification more generally.} The identification of the
contest parameters $(\eta, \alpha_f, \theta_f)$ relies on the joint
distribution of Epoch-measured model quality, compute build-out, and
observed investment flows. The reputational sensitivity $\gamma_f$ is
identified from the within-firm correlation of procurement-mode choice
with announced sustainability commitments and from variation across firms
in the implied willingness-to-pay for clean procurement.

\section{Physical Experiments}
\label{sec:experiments}

This section describes the physical experimental platform under
development at Stony Brook University's CEWIT/AERTC facility and its
integration with the structural model. The experimental platform produces
the variation in technical parameters that calibrates the $\Lambda$ map of
(\ref{eq:zeta_map}) and decomposes the observational ancillary-services
estimates of (\ref{eq:aux_did}).

\subsection{Platform Description}

The experimental platform consists of three coupled components: (i) a
GPU cluster operating as a ``micro data center'' (4 NVIDIA RTX
4090-class GPUs, 16-core CPU, 64 GB RAM, 2 TB SSD) running real AI
training and inference workloads; (ii) a microgrid platform at AERTC
providing controlled grid interconnection with variable generation,
storage, and load configurations; (iii) instrumentation at the
data-center-to-grid interface measuring power, voltage, current,
frequency, harmonic content, and reactive power at high temporal
resolution. A coupled digital twin developed in collaboration with
CEWIT integrates the physical measurements with simulation models of
the wider grid, enabling extrapolation from the micro-scale experiments
to grid-relevant magnitudes.

\paragraph{Why a coupled experimental--econometric platform.}
Observational data on AI grid impacts is fundamentally limited along
the dimensions of greatest policy interest. We observe zone-level outcomes
correlated with the presence of data centers, but we do not observe how a
specific change in chip generation, model precision, batch scheduling, or
power-electronics topology affects grid signatures. Conversely, pure
simulation of grid effects depends on parameters that are not well
identified from observational data alone. The coupled platform produces
the identifying variation in technical parameters that observational
data cannot.

\subsection{Three Tiers of Experiments}

Following the design specification, experiments are organized into three
tiers.

\paragraph{Tier 1: AI workload design.} We vary, on the micro data center,
parameters under software and chip-level control: quantization (FP32,
FP16, BF16, INT8, INT4), pruning ratio, model distillation, parallelization
(data parallel, tensor parallel, pipeline parallel), inter-GPU communication
overhead, training versus inference workload mix, batch size, precision
level, and scheduling pattern. Each configuration generates a measured load
profile at the data-center-to-grid interface.

\paragraph{Tier 2: Grid interconnection and power systems.} We vary the
interconnection design: power-electronics topology (converter type,
transformer configuration, power factor correction), load shaping and
demand-response strategy, backup configuration (BESS, hydrogen, diesel),
and the generation source (on-site solar, on-site wind, off-site PPA,
colocated). Each configuration generates measured signatures of power
quality, harmonic distortion, ramp rate, and inertia contribution at the
grid interconnect.

\paragraph{Tier 3: Integration with econometric models.} This tier is the
handoff, and it is worth being explicit about what crosses it. The experiments
deliver one object to the econometric model: a fitted shape map $\hat\Lambda$
(eq.~\ref{eq:zeta_map}) relating technical parameters to per-MW load
signatures, together with the engineering relationship $g_j(\cdot)$ from
signatures to per-MW ancillary-services demand. These are the scale-free
factors in the composition (\ref{eq:aux_composition}); the level
($K_{f,z,t}$, the fleet mix, the aggregation factor $a_{j,z}$) is supplied by
the observational data of \S\ref{sec:data}, not by the experiments. With
$\hat\Lambda$ and $g_j$ in hand the structural model can then answer the
counterfactual question the reduced form cannot: if a specific technical
change---INT8 inference, a power-factor-correction standard, coordinated batch
scheduling---were adopted across the fleet, the change perturbs
$\theta^{\text{tech}}$, $\hat\Lambda$ maps it to a new per-MW signature,
(\ref{eq:aux_composition}) re-aggregates it to a zone-level demand shift, and
the equilibrium re-solves for wholesale prices, reserve requirements, REC
effectiveness, reliability, and emissions. The experiments thus identify a
derivative that is unavailable in any observational dataset; the model turns
that derivative into a general-equilibrium counterfactual.

% \subsection{Contingency for Microgrid Access}

% The platform has two components---a physical microgrid at AERTC and a coupled
% digital twin---and the paper's identification does not stand or fall on the
% physical microgrid being online on any particular date. It is useful to state
% precisely what each component contributes and what is lost if the physical
% microgrid is delayed.

% The digital twin is sufficient to identify the \emph{Tier-1} shape map: the
% dependence of per-MW signatures on \emph{software- and chip-level} parameters
% (quantization, batch size, parallelization, workload mix) is governed by the
% GPUs and their power-delivery electronics, which the twin models directly and
% which the physical micro data center measures at the interface. What the
% \emph{physical microgrid} adds is the \emph{Tier-2} margin: the interaction of
% the load with real power-electronics topologies, transformer configurations,
% storage inverters, and on-site generation under controlled grid conditions,
% including transient and resonance behavior that a simulation can mis-state.

% If the physical microgrid is unavailable on the required timeline, we proceed
% on three legs and are explicit about the cost. First, the Tier-1 map is
% recovered from the micro data center and the twin and is unaffected. Second, the
% Tier-2 interactions are taken from the twin with their simulation uncertainty
% propagated forward, rather than asserted to be negligible; counterfactuals that
% turn on Tier-2 physics (e.g., a power-factor-correction standard) are reported
% with that uncertainty attached. Third---and this is the external check that
% keeps the contingency honest---the aggregate prediction is validated against the
% utility-level Whisker Labs THD distribution (\S\ref{sec:data}) and against the
% DiD-implied demand shift (\S\ref{sec:aux_market}); if the twin-only path cannot
% reproduce those aggregates, the DiD remains the binding number and the
% technical-decomposition counterfactuals are flagged as provisional pending
% physical measurement. The physical microgrid sharpens Tier-2 and removes the
% simulation uncertainty; it is not load-bearing for the Tier-1 results or for the
% aggregate impacts that the observational design already identifies.


\section{Results}
\label{sec:results}

% [To be completed after estimation. The results section will report:
% (i) estimates of the ancillary-services-impact coefficients
% $\hat\beta^j_k$ from the stacked DiD (\ref{eq:aux_did}) for each product
% and ISO; (ii) the experimentally-estimated $\hat\Lambda$ map and its
% cross-validation against the DiD-implied $\widehat{\Delta D^{\text{aux}}}$;
% (iii) recovered hyperscaler structural parameters
% $(\gamma_f, \alpha_f, \theta_f, \eta)$ with confidence sets; (iv) generator
% entry cost and financing-rate functions $\hat r_l(\sigma^2)$; (v) the
% implied PPA bargaining weights $\hat\zeta_f$.]

\subsection{Estimated Coefficients Short Run}

% [To be completed. Tables will report stacked-DiD coefficients for
% ancillary-services price by product (regulation up/down, synchronized
% reserve, primary reserve, voltage support) and by ISO (PJM, ERCOT,
% CAISO, MISO). Each estimate will be accompanied by pre-trend tests and
% the alternative-estimator robustness checks described in
% \S\ref{sec:est_short}.]

\section{Model Fit and Robustness}
\label{sec:fit}

% [To be completed. The model-fit exercises will include: (i) within-sample
% fit of observed PPA prices to model-predicted Nash-bargaining prices;
% (ii) within-sample fit of observed data-center capacity choices to model
% solutions; (iii) overidentifying tests using moments not used in
% estimation (e.g., observed REC prices, observed merchant procurement
% shares, observed generator entry by fuel type). Robustness exercises
% include alternative block partitions, alternative supply-curve
% specifications for the energy market, and alternative attribution rules
% for ancillary-services cost.]

\section{Impacts of AI on the Electricity System}
\label{sec:impacts}

This section uses the estimated structural model to decompose the impact
of AI demand growth on equilibrium outcomes in the electricity system,
holding policy fixed at its observed configuration. The decomposition
isolates the contribution of AI demand from other contemporaneous
factors (gas prices, weather, retirement of coal, growth in renewable
penetration) that confound reduced-form estimates.

\subsection{Impacts of AI on the Price of Electricity}

% [To be completed. The price-impact analysis will combine the
% DiD-estimated demand shift with the model-derived equilibrium response
% of the supply stack, including generator entry. The estimated price
% impact accounts for both the short-run rationing effect (which the
% reduced-form estimate captures) and the medium-run entry response
% (which the reduced-form estimate cannot capture but the structural
% model can simulate).]

\subsection{Impacts of AI on Emissions}

% [To be completed. The emissions analysis aggregates the firm-level
% $E_{f,t}$ from (\ref{eq:emissions}) across hyperscalers, decomposing
% the total into (i) direct emissions from increased fossil dispatch,
% (ii) avoided emissions from REC retirement under the granularity
% standard, (iii) emissions from backup running, and (iv) indirect
% emissions from changes in the generator entry mix. The decomposition
% isolates the effect of REC and PPA procurement choices on actual
% emissions, distinct from the accounting effect.]

\subsection{Impacts of AI on Power Quality}

% [To be completed. The power-quality analysis uses the estimated $\Lambda$
% map to translate observed AI load growth into predicted aggregate THD,
% voltage variability, and harmonic content at the zone level. The
% predictions are validated against Whisker Labs THD observations.]

\subsection{Impacts of AI on Ancillary Services Demand and Pricing}

% [To be completed. This is the central new descriptive result of the
% paper: the share of ancillary-services demand growth attributable to
% AI workloads, by product and ISO. The structural model also produces
% the implied increase in equilibrium ancillary-services prices and the
% distributional incidence of ancillary-services cost across customer
% classes.]

\section{Counterfactual Analysis}
\label{sec:cf}

The counterfactual exercises use the estimated structural model to
simulate equilibrium outcomes under modified choice environments. Each
counterfactual is implemented by re-solving the Markov-perfect equilibrium
with the modified constraint set, holding structural parameters fixed at
their estimated values, and comparing equilibrium portfolios, prices,
emissions, ancillary-services demand, and reliability to the baseline.
For each counterfactual we report both the partial-equilibrium effect
(holding generator entry fixed) and the general-equilibrium effect
(allowing entry to respond), since the difference is itself informative
about the channels through which the policy operates.

\subsection{Relocating Data Centers}

\paragraph{Locational redistribution of inference and training.}
The first counterfactual changes the geographic distribution of
hyperscaler workloads. Specifically, we consider three variations:
(i) shifting inference workloads from the current concentration zones
(Northern Virginia, central Ohio, Phoenix, Quincy/Boardman) to upstream
zones with greater renewable headroom; (ii) shifting training workloads
to zones with the lowest marginal carbon intensity; (iii) co-optimizing
training and inference allocation under a fleet-wide constraint on
ancillary-services costs.

The mechanism: shifting inference toward upstream zones reduces the
ancillary-services impact at congested downstream zones (because both
ramping and harmonic content of inference load are
proportional to the magnitude of the load locally) but increases
transmission utilization. The net welfare effect depends on the
elasticity of ancillary-services supply at each zone and on the
shadow price of transmission, both of which the structural model
recovers. The result is a counterfactual map of where AI workloads
``should'' locate to minimize a specified objective, and the magnitudes
of the implied benefits.

\subsection{Technical Changes at the Model and Data Center Level}

The second counterfactual uses the digital-twin-calibrated $\Lambda$ map
to evaluate the equilibrium effect of technical changes. Specific variations:
(i) fleet-wide adoption of INT8 inference quantization; (ii) widespread
deployment of distilled smaller inference models; (iii) standardized power-factor
correction at the data-center-to-grid interface; (iv) batch-scheduling
coordination across firms to smooth fleet-wide load. Each technical
change shifts $\vec\zeta$ for the affected workloads, which propagates
through (\ref{eq:aux_demand}) to ancillary-services demand and prices.
The counterfactual quantifies the total system benefit (lower
ancillary-services cost, improved power quality, reduced reserve
requirements) and the implied private incentive of the hyperscaler
to adopt the change.

This counterfactual is uniquely enabled by the experimental platform:
neither the reduced-form approach nor any purely observational structural
model can evaluate it, because the variation in technical parameters
required for identification is not present in observational data.

\subsection{Investment in Transmission or Generation Mandate}

The third counterfactual evaluates supply-side policy responses. Variations:
(i) targeted transmission expansion at congested arcs in the network,
implemented as an increase in $\bar T_{a,t}$ for selected arcs;
(ii) accelerated generator entry at specific zones via a build-rate
subsidy that reduces $I^{\text{build}}_{l,z,t}$; (iii) mandatory
generator entry tied to data-center build approvals (``bring-your-own-clean,''
treated separately below).

\subsection{Bring-Your-Own-Generation Mandate}

The fourth counterfactual restricts the procurement-mode choice set:
all incremental capacity at zone $z$ must be PPA or Colocation, forbidding
merchant and REC modes. Formally this imposes
$\sigma_{f,z,t} \in \{\text{PPA, Colo}\}$ for all $f$ at constrained zones.
The constraint interacts strongly with the clean-MW pool (\ref{eq:cleanpool}):
when $Z_{z,t}$ is small relative to data-center demand, the constraint
generates a wedge between hyperscaler willingness to pay and PPA price
that induces generator entry through the financing-channel mechanism
(\ref{eq:entry_ordering}). The counterfactual quantifies the magnitude
of the entry response, the implied equilibrium PPA price, and the
allocation effect across hyperscalers (firms with higher $\gamma_f$ find
the mandate cheaper to comply with because they would have chosen PPA
or Colo anyway).

\subsection{Elimination of the REC Market}

The fifth counterfactual removes the REC market entirely, setting
$\phi = 0$ for all firms (no REC retirements possible). This isolates
the operational effect of REC purchases: in the model, RECs reduce
attributable emissions but do not change the dispatched generation in
the short run. Eliminating RECs therefore changes attributable emissions
in (\ref{eq:emissions}) but not actual emissions in the short run.
In the long run, REC elimination removes one revenue stream for
renewable entrants, slowing entry and raising the actual marginal carbon
intensity of the grid. The counterfactual quantifies the relative
magnitudes of the accounting and operational effects.

\subsection{National REC Market}

The sixth counterfactual is the opposite of geographic REC restrictions:
$S^{\text{qual}}_{z,b,t} = \sum_{z'} S_{z',b,t}$ for all $z$ (RECs from
any zone qualify everywhere). This expands the effective clean-MW pool
seen by RECs but does not change (\ref{eq:cleanpool}) for PPA/Colo. The
exercise quantifies the cost-saving from broader REC eligibility, the
implied distributional effect (zones with abundant cheap renewables
become net REC exporters), and the equilibrium response in actual
emissions (which depends on whether REC eligibility expansion induces
entry or reshuffles existing supply).

\subsection{24/7 Carbon-free Electricity Mandate}

The seventh counterfactual sets the granularity standard to $h = 1$
(hourly matching), forcing each hyperscaler to retire RECs in the same
block as its load. The mechanism: hourly matching binds (\ref{eq:gran_req})
on a per-block basis, which redistributes value sharply toward
colocation-with-storage (because that combination is what shrinks the
timing wedge in the operational subgame). The counterfactual quantifies
the equilibrium shift toward colocation-plus-storage, the implied increase
in renewable-plus-storage entry, and the change in actual emissions.

\subsection{Changes in Ancillary Services Market Structure}

The eighth counterfactual modifies the ancillary-services market itself.
Variations: (i) introduction of a fast-frequency-response product
specifically designed for the ramping characteristics of AI load;
(ii) shift from proportional to incremental cost attribution in
(\ref{eq:hyper_aux_cost}), charging the marginal contributor at the
marginal price; (iii) a dedicated ancillary-services demand-response
product paying data centers to provide ramping and harmonic mitigation.
Each variation changes the marginal cost faced by the hyperscaler from
its ancillary-services contribution, which propagates back into the
technical-parameter choices via the dual-use storage channel
(\S\ref{sec:storage_dual}).

\subsection{Changing Model Designs Fleet-Wide}

The ninth counterfactual combines the technical-change counterfactual
with a coordination requirement: hyperscalers are required to disclose
load-shape signatures and coordinate scheduling. This is a hybrid of
the technical counterfactual and a market-design counterfactual.
The exercise quantifies the surplus available from coordinated scheduling
and the implied policy mechanism (mandatory disclosure, voluntary
coordination with regulatory backstop, etc.) required to achieve it.

\subsection{Composite Policy Scenarios}

We also report two composite scenarios that combine elements of the
single-lever counterfactuals: (i) a ``high-clean'' scenario combining
24/7 matching with BYOC and accelerated transmission expansion; and
(ii) a ``high-efficiency'' scenario combining technical-change mandates
with ancillary-services market redesign. The composite scenarios reveal
the policy complementarities (and substitutions) across levers, which
single-lever exercises cannot.

\section{Conclusions}
\label{sec:conclusion}

[To be completed after results.]

\section*{Appendix}
\label{sec:appendix}

\subsection*{A. Computational Implementation}

[To be completed. Implementation details for the moment-inequalities
estimator, the MPE solution algorithm, the digital-twin integration, and
the simulation of counterfactual equilibria. Code and replication
materials will be made available.]

\subsection*{B. Additional Derivations}

[To be completed. Full derivation of the procurement-mode cost
functions (\ref{eq:c_merch})--(\ref{eq:c_colo}), the REC market shadow
prices, and the proof that conditional separability (\ref{eq:separability})
holds under the maintained assumptions.]

\subsection*{C. Additional Robustness Exercises}

[To be completed.]

\newpage
\bibliographystyle{apalike}
\bibliography{references}

\end{document}
